\section{PROJETO ORIENTADO A OBJETOS}

\subsection{Aplicação de Princípios da Orientação a Objetos}

O projeto do \textbf{MedAgenda} fundamenta-se solidamente no paradigma Orientado a Objetos (POO), concretizando os pilares da programação modular na estruturação das entidades apresentadas no Diagrama de Classes (Seção 3.2) \cite{meyer1997}:

\subsubsection{1. Encapsulamento}
O encapsulamento foi sistematicamente aplicado em todas as entidades e módulos do sistema para proteger a integridade interna dos objetos contra acessos indevidos ou mutações inválidas de estado \cite{larman2004}:
\begin{itemize}
    \item Todos os atributos de estado e dados sensíveis (tais como \texttt{cpf}, \texttt{senhaHash}, \texttt{anamnese} e \texttt{diagnostico}) foram estritamente declarados com visibilidade privada (\texttt{-}).
    \item A leitura e mutação controlada de estado realizam-se apenas através de operações públicas semânticas de negócio (\texttt{+}) que aplicam validações de domínio. Por exemplo, a classe \texttt{Consulta} não expõe um método genérico \texttt{setStatus()}; em vez disso, provê métodos semânticos como \texttt{confirmar()}, \texttt{cancelar(motivo)} e \texttt{concluir()}, garantindo que as transições transacionem estritamente as regras modeladas no Diagrama de Estados \cite{meyer1997}.
\end{itemize}

\subsubsection{2. Herança}
A hierarquia de generalização e especialização foi explorada para modelar de forma limpa os atores e usuários corporativos da clínica:
\begin{itemize}
    \item A classe supertipo abstrata \textbf{\texttt{Usuario}} centraliza as propriedades e comportamentos comuns a qualquer pessoa autenticada na plataforma (\texttt{id}, \texttt{nome}, \texttt{email}, \texttt{senhaHash}, além das operações globais de sessão \texttt{login()} e \texttt{logout()}).
    \item As subclasses concretas \textbf{\texttt{Paciente}}, \textbf{\texttt{Medico}} e \textbf{\texttt{Administrador}} herdam as propriedades fundamentais de \texttt{Usuario} e adicionam atributos específicos das suas obrigações clínicas e corporativas (como \texttt{crm} e \texttt{especialidade} para Médicos; \texttt{cpf} e \texttt{dataNascimento} para Pacientes; \texttt{nivelAcesso} para Administradores).
\end{itemize}

\subsubsection{3. Polimorfismo}
O polimorfismo de inclusão e sobrescrita de métodos foi aplicado no roteamento dinâmico e na verificação de permissões de acesso:
\begin{itemize}
    \item O método abstrato \texttt{obterPainelInicial()} definido em \texttt{Usuario} é polimorficamente implementado em cada subclasse específica, retornando ao navegador do cliente o painel ou dashboard de navegação personalizado e ergonômico correspondente a cada perfil de acesso \cite{larman2004}.
    \item O mecanismo central de controle de permissão processa referências polimórficas do tipo genérico \texttt{Usuario}, resolvendo dinamicamente em tempo de execução as regras lógicas específicas de cada subclasse concreta.
\end{itemize}

\subsection{Justificativa de Decisões de Projeto e Relacionamentos}

\begin{itemize}
    \item \textbf{Por que a classe \texttt{HorarioDisponivel} foi modelada de forma autônoma e separada de \texttt{Medico}?} \newline
    Em vez de armazenar horários em um array de strings simples dentro da entidade Médico, criou-se a entidade relacional autônoma \texttt{HorarioDisponivel}. Essa decisão arquitetural viabiliza consultas SQL altamente otimizadas e filtradas (por exemplo: listar em milissegundos todos os horários livres para a especialidade Cardiologia na sexta-feira), evita concorrência e bloqueios transacionais durante agendamentos simultâneos e assegura que a entidade \texttt{Medico} permanezca com alta coesão estrutural \cite{larman2004, silberschatz2020}.
    \item \textbf{Por que foi escolhida a relação de Composição forte entre \texttt{Consulta} e \texttt{Prontuario}?} \newline
    O prontuário eletrônico constitui o registro clínico e legal indissociável de uma consulta médica finalizada. Ele não possui razão de existir no banco de dados sem estar subordinado a uma consulta específica que lhe originou. A composição forte assegura que a integridade referencial e o ciclo de vida transacional do prontuário permaneçam vinculados e dependentes da consulta \cite{booch2005}.
\end{itemize}

\subsection{Distribuição de Responsabilidades entre as Classes}

A distribuição de responsabilidades e métodos seguiu rigorosamente os padrões GRASP (\textit{General Responsibility Assignment Software Patterns}) e os princípios de design SOLID \cite{larman2004, martin2019}:
\begin{itemize}
    \item \textbf{Especialista na Informação (\textit{Information Expert})}: A classe \texttt{Consulta} foi designada como o especialista natural responsável por validar se o cancelamento pode ser acionado pelo paciente, pois é ela que detém a informação exata da propriedade \texttt{dataHora} da consulta e conhece a regra lógicocronológica das 24 horas de antecedência \cite{larman2004}.
    \item \textbf{Controlador (\textit{Controller})}: As classes de controlador da camada de aplicação (como \texttt{ConsultaController} e \texttt{AutenticacaoController}) atuam como intermediárias transacionais entre as requisições HTTP da interface web e a lógica de negócios da clínica, orquestrando fluxos sem acumular regras de domínio que pertencem às entidades \cite{fowler2002}.
    \item \textbf{Princípio da Responsabilidade Única (\textit{Single Responsibility Principle - SRP})}: O registro de anamnese, prescrição terapêutica e diagnóstico com CID-10 foi isolado na entidade própria \texttt{Prontuario}, evitando sobrecarregar a classe \texttt{Consulta} com campos clínicos longos, preservando a coesão do modelo e facilitando o controle de sigilo médico \cite{martin2019}.
\end{itemize}
